% Book-specific commands, environments and metadata.

% Title metadata used by the frontmatter pages.
\newcommand{\thebooktitle}{Độ trung thực trong XAI}
\newcommand{\thebooksubtitle}{Từ LIME đến chuỗi suy luận}
\newcommand{\thebookauthor}{Giang Dang}

% Inline code identifiers (method names, hyperparameters, library symbols).
% This book ships no code that runs; \code marks a name, never a program.
\newcommand{\code}[1]{\texttt{#1}}

% ---------------------------------------------------------------------------
% Glossed terms
% ---------------------------------------------------------------------------
% This book is written in Vietnamese for readers who will go on to read the
% 32 corpus papers in English. A term the book translates carries its English
% original in parentheses, on the cadence the SPEC's gloss regime sets: a term
% a chapter owns is glossed once per section, a borrowed term once per chapter.
%
%   \tn{độ trung thực}{faithfulness}  ->  độ trung thực (faithfulness)
%
% Terms the book keeps in English (attribution, baseline, benchmark and
% friends) are typed plainly and take no parentheses. Which terms fall on
% which side is settled once in appendix B and does not drift; the Gloss check
% in check-chapter.psd1 reads that appendix and enforces the cadence.
\newcommand{\tn}[2]{#1 (#2)}

% ---------------------------------------------------------------------------
% Theorem environments
% ---------------------------------------------------------------------------
% One shared counter, numbered by chapter, so a later chapter can cite
% "định nghĩa 1.2" with no ambiguity about which sequence it belongs to.
% Babel's Vietnamese locale already supplies \proofname as "Chứng minh".
\theoremstyle{definition}
\newtheorem{definition}{Định nghĩa}[chapter]
\newtheorem{example}[definition]{Ví dụ}

\theoremstyle{plain}
\newtheorem{theorem}[definition]{Định lý}
\newtheorem{lemma}[definition]{Bổ đề}
\newtheorem{proposition}[definition]{Mệnh đề}
\newtheorem{corollary}[definition]{Hệ quả}

\theoremstyle{remark}
\newtheorem{remark}[definition]{Nhận xét}

% The algorithm float sets its own caption name rather than taking one from
% babel, so left alone it prints "Algorithm 1" in the middle of a Vietnamese
% page. Every pseudo-code listing in this book is one of these floats.
\floatname{algorithm}{Thuật toán}

% ---------------------------------------------------------------------------
% Chapter apparatus
% ---------------------------------------------------------------------------
% Grey and rules only, no colour: these books are likely to be read in mono.

% End-of-chapter summary box: 3 to 6 bullets restating the chapter's claims
% and definitions, nothing that did not already appear in the body.
\newtcolorbox{tomtat}{
  breakable, enhanced, sharp corners,
  colback = black!3, colframe = black!55,
  boxrule = 0.4pt, left = 3mm, right = 3mm, top = 2mm, bottom = 2mm,
  title = {Tóm tắt},
  fonttitle = \bfseries\small,
  coltitle = black, colbacktitle = black!12,
}

% Review questions follow the summary box. Starred heading: nothing
% cross-references a question, and eighteen chapters of them would otherwise
% flood the table of contents.
\newcommand{\cauhoi}{\section*{Câu hỏi ôn tập}}

\newlist{cauhoids}{enumerate}{1}
\setlist[cauhoids]{label = \arabic*., leftmargin = *, itemsep = 0.4em,
  topsep = 0.4em}
